\chapter{}
\label{chap:3}

\section{Умова}

Знайти коефіцієнти Фур’є булевої функції $f$, якщо її вектор значень дорівнює:
\begin{enumerate}[label=\alph*)]
    \item (1, 1, 0, 1, 0, 1, 1, 0);
    \item (1, 1, 1, 0, 0, 1, 0, 1).
\end{enumerate}

\section{Розв'язання}

\indent Спершу я для кращого розуміння розписав швидке перетворення Адамара, яке ми розбирали на лекції і застосовували на 
практиці. Коефіцієнти Фур'є обчислюються за формулою $C_{f} = H_{n} \cdot f$, де $H_n$ -- матриця Адамара порядку $2^n$.

Замість прямого множення (що потребує $2^{2n}$ операцій) використовуємо рекурентне співвідношення для матриць 
Адамара, яке доводили в попередній задачі.

Для вектора $a = \begin{pmatrix} a_0 \\ a_1 \end{pmatrix}$, де через $a_0, a_1$ -- позначено верхню та нижню половину 
відповідно, маємо:
\begin{equation*}
    H_n a = \begin{pmatrix} H_{n-1}(a_0 + a_1) \\ H_{n-1}(a_0 - a_1) \end{pmatrix}.
\end{equation*}
У нас $n = 3$, тому алгоритм виконує усього $3$ кроки. Працюємо за правилом: розбиваємо вектор на дві половини. 
Верхню та нижню половини покомпонентно додаємо (для верхньої) та віднімаємо (для нижньої), і так доки не розіб'ємо 
по два елементи. Потім робимо перевірку рівності Парсеваля для коефіцієнтів Фур'є:
\begin{equation*}
    \frac{1}{2^n}\sum\limits_{\alpha \in V_n} C_f(\alpha)^2 = \sum\limits_{x \in V_n} f(x)^2 = \operatorname{wt}(f).
\end{equation*}

\newpage
\subsection*{а) $f = (1, 1, 0, 1, 0, 1, 1, 0)$}

\begin{table}[htpb]
    \centering
    \begin{tblr}{
        colspec = {c | c | c | c},
        row{1} = {font=\bfseries},
        hline{2,10} = {solid}, % pivot lines
        hline{6} = {1}{solid}, % first column
        hline{4,6,8} = {2}{solid}, % second column
        hline{2,3,4,5,6,7,8,9} = {3}{solid} % third column
    }
        $f$ & Крок 1 & Крок 2 & $C_f$ \\
        1   & 1      & 2      & 5     \\
        1   & 2      & 3      & $-1$  \\
        0   & 1      & 0      & 1     \\
        1   & 1      & 1      & $-1$  \\
        0   & 1      & 0      & 1     \\
        1   & 0      & 1      & $-1$  \\
        1   & $-1$   & 2      & 1     \\
        0   & 1      & $-1$   & 3     \\
    \end{tblr}
    \caption{Швидке перетворення Адамара для $f = (1,1,0,1,0,1,1,0)$}
\end{table}

\noindent У нас $\operatorname{wt}(f) = 5$. Робимо перевірку:
\begin{equation*}
    \frac{5^2 + ({-}1)^2 + 1^2 + ({-}1)^2 + 1^2 + ({-}1)^2 + 1^2 + 3^2}{2^3}
    = \frac{40}{8} = 5 = \operatorname{wt}(f).
\end{equation*}

\noindent Відповідь: $C_f = (5,\; {-}1,\; 1,\; {-}1,\; 1,\; {-}1,\; 1,\; 3)$

\newpage
\subsection*{б) $f = (1, 1, 1, 0, 0, 1, 0, 1)$}

\begin{table}[htpb]
    \centering
    \begin{tblr}{
        colspec = {c | c | c | c},
        row{1} = {font=\bfseries},
        hline{2,10} = {solid}, % pivot lines
        hline{6} = {1}{solid}, % first column
        hline{4,6,8} = {2}{solid}, % second column
        hline{2,3,4,5,6,7,8,9} = {3}{solid} % third column
    }
        $f$ & Крок 1 & Крок 2 & $C_f$ \\
        1   & 1      & 2      & 5     \\
        1   & 2      & 3      & $-1$  \\
        1   & 1      & 0      & 1     \\
        0   & 1      & 1      & $-1$  \\
        0   & 1      & 2      & 1     \\
        1   & 0      & $-1$   & 3     \\
        0   & 1      & 0      & 1     \\
        1   & $-1$   & 1      & $-1$  \\
    \end{tblr}
    \caption{Швидке перетворення Адамара для $f = (1,1,1,0,0,1,0,1)$}
\end{table}

\noindent Теж $\operatorname{wt}(f) = 5$. Перевіряємо:
\begin{equation*}
    \frac{5^2 + ({-}1)^2 + 1^2 + ({-}1)^2 + 1^2 + 3^2 + 1^2 + ({-}1)^2}{2^3}
    = \frac{40}{8} = 5 = \operatorname{wt}(f).
\end{equation*}

\noindent Відповідь: $C_f = (5,\; {-}1,\; 1,\; {-}1,\; 1,\; 3,\; 1,\; {-}1)$
